\begin{solucion}
Sea $S_n$ el conjunto de todas las permutaciones de $[n]=\{1,2,\dots,n\}$ y, para cada
$1\le i\le n$, denotemos  
\[
A_i \;=\;\bigl\{\sigma\in S_n \mid \sigma(i)=i\bigr\},
\]
es decir, el subconjunto de permutaciones que fijan al elemento $i$.  
El número de desarreglos buscado es el cardinal de la intersección de los
complementos $\overline{A_i}=S_n\setminus A_i$:
\[
D_n \;=\;\Bigl|\bigcap_{i=1}^{n}\overline{A_i}\Bigr|.
\]

Aplicamos ahora el \emph{Principio de Inclusión–Exclusión}: para cualquier colección
finita de subconjuntos $A_1,\dots,A_n$ de un conjunto $S_n$ se tiene cumple que:
\begin{align*}
    \Bigl|\bigcap_{i=1}^{n}\overline{A_i}\Bigr|
\;&=\;
|S_n|\;-\!\!\sum_{j=1}^{n}(-1)^{\,j-1}
\!\!\!\sum_{\{i_1,\dots,i_j\}}\!\!\!
\bigl|A_{i_1}\cap\dots\cap A_{i_j}\bigr|\\
D_n \;&=\; n!
-\sum_{j=1}^{n}(-1)^{\,j-1}
\!\!\!\sum_{\{i_1,\dots,i_j\}}\!\!\!
\bigl|A_{i_1}\cap\dots\cap A_{i_j}\bigr|.
\end{align*}


Fijado un $j\in\{1,\dots,n\}$ y un subconjunto $\{i_1,\dots,i_j\}\subset[n]$, con la restricción de que
las $j$ posiciones $i_1,\dots,i_j$ sean puntos fijos, lo cual nos deja 
libremente sobre los $n-j$ elementos restantes que se pueden permutar. Por lo tanto  
\[
\bigl|A_{i_1}\cap\dots\cap A_{i_j}\bigr|\;=\;(n-j)!,
\]
independientemente del subconjunto elegido. Como hay $\binom{n}{j}$ maneras de
seleccionar $j$ índices entonces:
\begin{align*}
    \sum_{\{i_1,\dots,i_j\}}\bigl|A_{i_1}\cap\dots\cap A_{i_j}\bigr| = \binom{n}{j}(n-j)!
\end{align*}
sustituimos en la expresión anterior y obtenemos
\begin{align*}
D_n\;&=\; n!-\sum_{j=1}^{n}(-1)^{\,j-1}\binom{n}{j}(n-j)!.
\end{align*}


Reagrupamos el término $j=0$ (igual a $n!$) en la suma:
\[
D_n \;=\;\sum_{j=0}^{n}(-1)^{\,j}\binom{n}{j}(n-j)!
\;=\;\sum_{j=0}^{n}(-1)^{\,j}\frac{n!}{j!}
\;=\; n!\sum_{j=0}^{n}\frac{(-1)^{\,j}}{j!}.
\]
\end{solucion}